% GENERATED FILE. Edit canonical lessons, then run npm run generate:books.
% canonical-source-hash: fnv1a64:198a40316bf77eab
% canonical-lessons: ML-C539-mayannuka, ML-C539-kitaykkuka, ML-C539-niruka, ML-C539-bhaksanakramam, ML-C539-posakaharam

\chapter{Dozing, Panting and Eating Well}
\label{ch:ml-dozing-panting-and-eating-well}
\hlchaptermodality{539}

\begin{chapteropening}
\textbf{By the end of this chapter:} \emph{I can say to doze off, to pant, to smart, a diet and nutritious food.}

Complete the last lesson of chapter 539: I can say to doze off, to pant, to smart, a diet and nutritious food.
\end{chapteropening}

\section[\texorpdfstring{\ml{മയങ്ങുക} (mayaṅṅuka)}{മയങ്ങുക (mayaṅṅuka)}]{\ml{മയങ്ങുക} (mayaṅṅuka) — to doze off; to drowse}

\label{lesson:ML-C539-mayannuka}



\begin{quote}
Before the new one: say the Malayalam for to hurt, then the Malayalam for to pass out.
\end{quote}

\subsection*{You'll want to know: \ml{മയങ്ങുക}}
\textbf{\ml{മയങ്ങുക}} — \emph{mayaṅṅuka} — \textquotedblleft{}to doze off; to drowse\textquotedblright{}.

\begin{quote}
\ml{ഊണ്} \ml{കഴിഞ്ഞ്} \ml{മുത്തശ്ശി} \ml{കസേരയിൽ} \ml{ഇരുന്ന്} \ml{മയങ്ങി}.
\end{quote}

\begin{quote}
\emph{ūṇŭ kaḻiññŭ muttaśśi kasērayil irunnŭ mayaṅṅi.}
\end{quote}

\begin{quote}
After lunch grandmother dozed off in her chair.
\end{quote}

Past: \ml{മയങ്ങി}, \emph{mayaṅṅi}.

A light sleep; the noun is \ml{മയക്കം}, drowsiness, as after a meal or a medicine.

\begin{cousinweb}
\textbf{\ml{മയങ്ങുക}} (\emph{mayaṅṅuka}) is native Dravidian: DEDR 4706 gives it as to be drowsy or giddy, with \ml{മയക്കം}, drowsiness (DEDR 4706).
\end{cousinweb}

\begin{grammarlens}[title={What you've built}]
One more word for describing what you lived and felt.
\end{grammarlens}

\subsection*{Your turn}
\begin{itemize}
\raggedright
  \item \emph{Say it:} \emph{mayaṅṅuka}
  \item \emph{Say it:} \emph{mayaṅṅuka}, once more
  \item \emph{Say it:} \emph{ūṇŭ kaḻiññŭ muttaśśi kasērayil irunnŭ mayaṅṅi.} — After lunch grandmother dozed off in her chair.
  \item \emph{Recall:} say the Malayalam for to hurt, then the Malayalam for to pass out, then say \emph{mayaṅṅuka} again
\end{itemize}

\begin{tcolorbox}[breakable,colback=teal!4,colframe=teal!35!black,title={Before you move on}]
What does \ml{മയങ്ങുക} mean? (\textquotedblleft{}To doze off; to drowse\textquotedblright{}.) Say it once more.
\end{tcolorbox}
\section[\texorpdfstring{\ml{കിതയ്ക്കുക} (kitaykkuka)}{കിതയ്ക്കുക (kitaykkuka)}]{\ml{കിതയ്ക്കുക} (kitaykkuka) — to pant; to gasp}

\label{lesson:ML-C539-kitaykkuka}



\begin{quote}
Before the new one: say the Malayalam for to pass out, then the Malayalam for to doze off.
\end{quote}

\subsection*{You'll want to know: \ml{കിതയ്ക്കുക}}
\textbf{\ml{കിതയ്ക്കുക}} — \emph{kitaykkuka} — \textquotedblleft{}to pant; to gasp\textquotedblright{}.

\begin{quote}
\ml{പടി} \ml{കയറി} \ml{വന്ന} \ml{മുത്തച്ഛൻ} \ml{കിതച്ചു}.
\end{quote}

\begin{quote}
\emph{paṭi kayaṟi vanna muttacchan kitaccu.}
\end{quote}

\begin{quote}
Grandfather was panting after climbing the steps.
\end{quote}

Past: \ml{കിതച്ചു}, \emph{kitaccu}.

The past is \ml{കിതച്ചു}; the noun is \ml{കിതപ്പ്}.

\begin{cousinweb}
\textbf{\ml{കിതയ്ക്കുക}} (\emph{kitaykkuka}) looks native, but DEDR has no entry that clearly holds it, so its origin is not settled.
\end{cousinweb}

\begin{grammarlens}[title={What you've built}]
One more word for describing what you lived and felt.
\end{grammarlens}

\subsection*{Your turn}
\begin{itemize}
\raggedright
  \item \emph{Say it:} \emph{kitaykkuka}
  \item \emph{Say it:} \emph{kitaykkuka}, once more
  \item \emph{Say it:} \emph{paṭi kayaṟi vanna muttacchan kitaccu.} — Grandfather was panting after climbing the steps.
  \item \emph{Recall:} say the Malayalam for to pass out, then the Malayalam for to doze off, then say \emph{kitaykkuka} again
\end{itemize}

\begin{tcolorbox}[breakable,colback=teal!4,colframe=teal!35!black,title={Before you move on}]
What does \ml{കിതയ്ക്കുക} mean? (\textquotedblleft{}To pant; to gasp\textquotedblright{}.) Say it once more.
\end{tcolorbox}
\section[\texorpdfstring{\ml{നീറുക} (nīṟuka)}{നീറുക (nīṟuka)}]{\ml{നീറുക} (nīṟuka) — to smart; to sting}

\label{lesson:ML-C539-niruka}



\begin{quote}
Before the new one: say the Malayalam for to doze off, then the Malayalam for to pant.
\end{quote}

\subsection*{You'll want to know: \ml{നീറുക}}
\textbf{\ml{നീറുക}} — \emph{nīṟuka} — \textquotedblleft{}to smart; to sting\textquotedblright{}.

\begin{quote}
\ml{മുറിവിൽ} \ml{മരുന്ന്} \ml{വച്ചപ്പോൾ} \ml{നീറി}.
\end{quote}

\begin{quote}
\emph{muṟivil marunnŭ vaccappōḷ nīṟi.}
\end{quote}

\begin{quote}
It stung when the medicine went on the wound.
\end{quote}

Past: \ml{നീറി}, \emph{nīṟi}.

The sting of a cut or of smoke in the eyes; a grief can also \ml{നീറുക} inside.

\begin{cousinweb}
\textbf{\ml{നീറുക}} (\emph{nīṟuka}) is native Dravidian: DEDR 3693 lists it with words for burning to ashes and for a grief that burns (DEDR 3693).
\end{cousinweb}

\begin{grammarlens}[title={What you've built}]
One more word for describing what you lived and felt.
\end{grammarlens}

\subsection*{Your turn}
\begin{itemize}
\raggedright
  \item \emph{Say it:} \emph{nīṟuka}
  \item \emph{Say it:} \emph{nīṟuka}, once more
  \item \emph{Say it:} \emph{muṟivil marunnŭ vaccappōḷ nīṟi.} — It stung when the medicine went on the wound.
  \item \emph{Recall:} say the Malayalam for to doze off, then the Malayalam for to pant, then say \emph{nīṟuka} again
\end{itemize}

\begin{tcolorbox}[breakable,colback=teal!4,colframe=teal!35!black,title={Before you move on}]
What does \ml{നീറുക} mean? (\textquotedblleft{}To smart; to sting\textquotedblright{}.) Say it once more.
\end{tcolorbox}
\section[\texorpdfstring{\ml{ഭക്ഷണക്രമം} (bhakṣaṇakramaṁ)}{ഭക്ഷണക്രമം (bhakṣaṇakramaṁ)}]{\ml{ഭക്ഷണക്രമം} (bhakṣaṇakramaṁ) — a diet; eating habits}

\label{lesson:ML-C539-bhaksanakramam}



\begin{quote}
Before the new one: say the Malayalam for to pant, then the Malayalam for to smart.
\end{quote}

\subsection*{You'll want to know: \ml{ഭക്ഷണക്രമം}}
\textbf{\ml{ഭക്ഷണക്രമം}} — \emph{bhakṣaṇakramaṁ} — \textquotedblleft{}a diet; eating habits\textquotedblright{}.

\begin{quote}
\ml{വൈദ്യൻ} \ml{പറഞ്ഞതുപോലെ} \ml{ഞാൻ} \ml{ഭക്ഷണക്രമം} \ml{മാറ്റി}.
\end{quote}

\begin{quote}
\emph{vaidyan paṟaññatupōle ñān bhakṣaṇakramaṁ māṟṟi.}
\end{quote}

\begin{quote}
I changed my diet as the doctor said.
\end{quote}

\ml{ഭക്ഷണം} food + \ml{ക്രമം} order, an arrangement.

\begin{cousinweb}
\textbf{\ml{ഭക്ഷണക്രമം}} (\emph{bhakṣaṇakramaṁ}) is a learned compound built from \ml{ഭക്ഷണം}, from Sanskrit \emph{bhakṣaṇa}, \textquotedblleft{}eating\textquotedblright{} (Monier-Williams), and \emph{krama}, \textquotedblleft{}a step, a series, an order\textquotedblright{} (CDIAL 3577).
\end{cousinweb}

\begin{grammarlens}[title={What you've built}]
One more word for describing what you lived and felt.
\end{grammarlens}

\subsection*{Your turn}
\begin{itemize}
\raggedright
  \item \emph{Say it:} \emph{bhakṣaṇakramaṁ}
  \item \emph{Say it:} \emph{bhakṣaṇakramaṁ}, once more
  \item \emph{Say it:} \emph{vaidyan paṟaññatupōle ñān bhakṣaṇakramaṁ māṟṟi.} — I changed my diet as the doctor said.
  \item \emph{Recall:} say the Malayalam for to pant, then the Malayalam for to smart, then say \emph{bhakṣaṇakramaṁ} again
\end{itemize}

\begin{tcolorbox}[breakable,colback=teal!4,colframe=teal!35!black,title={Before you move on}]
What does \ml{ഭക്ഷണക്രമം} mean? (\textquotedblleft{}A diet; eating habits\textquotedblright{}.) Say it once more.
\end{tcolorbox}
\section[\texorpdfstring{\ml{പോഷകാഹാരം} (pōṣakāhāraṁ)}{പോഷകാഹാരം (pōṣakāhāraṁ)}]{\ml{പോഷകാഹാരം} (pōṣakāhāraṁ) — nutritious food; nutrition}

\label{lesson:ML-C539-posakaharam}



\begin{quote}
Before the new one: say the Malayalam for to smart, then the Malayalam for a diet.
\end{quote}

\subsection*{You'll want to know: \ml{പോഷകാഹാരം}}
\textbf{\ml{പോഷകാഹാരം}} — \emph{pōṣakāhāraṁ} — \textquotedblleft{}nutritious food; nutrition\textquotedblright{}.

\begin{quote}
\ml{കുട്ടികൾക്ക്} \ml{പോഷകാഹാരം} \ml{കിട്ടണം}.
\end{quote}

\begin{quote}
\emph{kuṭṭikaḷkkŭ pōṣakāhāraṁ kiṭṭaṇaṁ.}
\end{quote}

\begin{quote}
Children must get nutritious food.
\end{quote}

\ml{പോഷകം} nourishing + \ml{ആഹാരം} food, joined into one long \ml{ആ}.

\begin{cousinweb}
\textbf{\ml{പോഷകാഹാരം}} (\emph{pōṣakāhāraṁ}) is a learned compound built from Sanskrit \emph{poṣaka}, \textquotedblleft{}nourishing\textquotedblright{} (Monier-Williams), and \emph{āhāra}, \textquotedblleft{}taking food, food\textquotedblright{} (CDIAL 1544).
\end{cousinweb}

\begin{grammarlens}[title={What you've built}]
One more word for describing what you lived and felt.
\end{grammarlens}

\subsection*{Your turn}
\begin{itemize}
\raggedright
  \item \emph{Say it:} \emph{pōṣakāhāraṁ}
  \item \emph{Say it:} \emph{pōṣakāhāraṁ}, once more
  \item \emph{Say it:} \emph{kuṭṭikaḷkkŭ pōṣakāhāraṁ kiṭṭaṇaṁ.} — Children must get nutritious food.
  \item \emph{Recall:} say the Malayalam for to smart, then the Malayalam for a diet, then say \emph{pōṣakāhāraṁ} again
\end{itemize}

\begin{tcolorbox}[breakable,colback=teal!4,colframe=teal!35!black,title={Before you move on}]
What does \ml{പോഷകാഹാരം} mean? (\textquotedblleft{}Nutritious food; nutrition\textquotedblright{}.) Say it once more.
\end{tcolorbox}
